\newpage
\hypertarget{para:postoincore}{}

\mypoemtitle{Io m’aggio posto in core a Dio servire - Giacomo da Lentini}{Sonetto}

\textbf{Commento introduttivo:}\\
In questo sonetto, Giacomo da Lentini affronta il rapporto tra amore profano e beatitudine celeste, teorizzando la necessità della presenza della donna anche nel regno di Dio per raggiungere la vera gioia.

\vspace{1.5em}

\begin{adjustwidth}{2em}{0pt}
\poemlines{5}
\verselinenumbersleft
\begin{verse}
Io m’aggio posto in core a Dio servire, \label{posto1}\\
com’io potesse gire in paradiso, \label{posto2}\\
al santo loco, c’aggio audito dire, \\
o’ si mantien sollazzo, gioco e riso. \label{posto4}\\!

     Sanza mia donna non vi voria gire, \\
quella c’à blonda testa e claro viso, \label{posto6}\\
che sanza lei non poteria gaudere, \\ 
estando da la mia donna diviso. \\!

     Ma no lo dico a tale intendimento, \\
perch’io pecato ci volesse fare; \\
se non veder lo suo bel portamento \\!

     e lo bel viso e ’l morbido sguardare: \\
che·l mi teria in gran consolamento, \\
veggendo la mia donna in ghiora stare. \label{posto14}\\
\end{verse}
\end{adjustwidth}

\vspace{1.5em}

\begin{commentaries}
    \scolio[\ref{posto2}]{gire} Andare.
    \scolio[\ref{posto6}]{blonda testa...} Canone di bellezza che diverrà standard nella lirica.
    \scolio[\ref{posto14}]{In ghiora stare} Stare nella gloria del paradiso. ghiora è un termine toscano sicuramente inserito dai copisti
\end{commentaries}

\vspace{0.5em}
{\color{Tomato!40}\hrule height 0.5pt}
\vspace{1em}

\subsubsection*{Analisi}
\begin{quote}
\small \itshape
A partire dalla prima quartina del sonetto, il poeta mette in luce la sua intenzione ovvero quella di servire ed avere fede in Dio per poter andare in paradiso e scoprire se è vero quello che gli è stato detto, ovvero che il paradiso sia un luogo felice e gioioso. già nella prima quartina, ed in particolar modo al vv.4, si crea una relazione tra la chiesa e la corte andando a vedere il paradiso come un ambiente di corte quasi giullaresco.

Nella seconda quartina viene introdotto un'ulteriore elemento presente nella poesia Sicliana, ovvero l'amore messo sullo stesso livello, se non su un livello superiore rispetto a quello della chiesa: Giacomo dice infatti di voler andare in paradiso solo a condizione che ci sia la la sua donna perché essendo da ella diviso non potrebbe essere felice nemmeno in paradiso

Nelle terzine il poeta mette le mani avanti, non vuole la donna per poter fare cose impure ma per le sue qualità spirituali e visive, in modo da poterla contemplare.
\end{quote}